\appendix
\setcounter{section}{1}
\section{Proof of \Cref{lemma:samp-opt}}\label{sec:proof-of-samp}
\begin{lemma}[\Cref{lemma:samp-opt}, restated. \cite{RVW01,Gol11}]
For every $\delta,\eps>0$ and integer $m$, there exists a $(\eps,\delta)$-sampler $f:\bits{n}\times\bits{d}\to\bits{m}$ s.t. $d=O(\log\log(1/\delta)+\log(1/\eps))$ and $n=m+O(\log(1/\delta))+O(\log(1/\eps))$. Moreover, for every $x,y$, $f(x,y)$ can be computed in space $O(m+\log(1/\delta)+\log(1/\eps))$.
\end{lemma}

We will use the equivalence between seeded randomness extractor and oblivious sampler by Zuckerman~\cite{Zuc97}. To achieve the parameter we need, we need a ``high-entropy seeded extractor" such that the seed length only depends on entropy loss but not the length of source. We will use the standard ``block-source" construction for high-entropy extractor which can be found in \cite{GW97,RVW01}. For simplicity, we will use simple composition instead of zig-zag composition~\cite{RVW01} because we are not aiming for optimal entropy loss. We will use the following standard lemmas for the extractor construction. Some of the following lemmas are implicit in their original source, and we recommend the readers to see \cite{GUV09,Vad12} for a proof.
\begin{definition}[\cite{CG88}]
$(X_1,X_2)$ is a \emph{$(k_1,k_2)$-block source} if $X_1$ is a $k_1$-source, and for every $x_1\in\Supp(X)$, $X_2$ conditioned on $X_1=x_1$  is a $k_2$-source.
\end{definition}
\begin{lemma}[\cite{GW97}]\label{lemma:block}
Let $X\in\bits{n}$ be a $(n-\Delta)$ source. Then for every integer $0\le t\le n $, $X$ is $\eps$-close to a $(t-\Delta,n-t-\Delta-\log(1/\eps))$-block source $(X_1,X_2)$ where $X_1\in\bits{t}$ and $X_2\in\bits{n-t}$.
\end{lemma}
\begin{lemma}[\cite{NZ96}]\label{lemma:compo}
Let $E_1:\bits{n_1}\times\bits{d}\to\bits{d_2}$ be a $(k_1,\eps_1)$ extractor and $E_2:\bits{n_2}\times\bits{d_2}\to\bits{m}$ be a $(k_2,\eps_2)$ extractor. Define $E((x_1,x_2),s)=E_1(x_2,E_2(x_1,s))$. Then for every $(k_1,k_2)$-block source $(X_1,X_2)\in\bits{n_1}\times\bits{n_2}$, $E((X_1,X_2),U_d)$ is ${(\eps_1+\eps_2)}$-close to uniform.
\end{lemma}
\begin{lemma}[\cite{GW97}]\label{lemma:GW}
For every $\eps,\Delta>0$ and integer $n$ there exists a $(n-\Delta,\eps)$ extractor $E:\bits{n}\times\bits{d}\to\bits{n}$ with $d=O(\Delta+\log(1/\eps))$, and for every $x,y$, $E(x,y)$ can be computed in space $O(n+\log(1/\eps))$.
\end{lemma}
\begin{lemma}[\cite{GUV09,KNW08}]\label{lemma:GUV}
For every $\eps>0$, integer $m>0$ and $n\ge 2m$, there exists a $(2m,\eps)$ extractor $E:\bits{n}\times\bits{d}\to\bits{m}$ with $d=O(\log m + \log(1/\eps))$, and for every $x,y$, $E(x,y)$  can be computed in space $O(m+\log(1/\eps))$.
\end{lemma}
\begin{lemma}[\cite{Zuc97}]\label{lemma:Zuc}
Every $(n-\log(1/\delta)-1,\eps)$-extractor is a $(\eps,\delta)$-sampler.
\end{lemma}
\noindent
Now we show how to construct the sampler we need, and that it is indeed space efficient.
\begin{proof}
Let $\Delta=\log(1/\delta)+1$. Let $E_1:\bits{m}\times\bits{d_1}\to\bits{m}$ be an $(m-\Delta,\eps/3)$-extractor from Lemma~\ref{lemma:GW}, w.l.o.g. assume that $d_1\ge \Delta+\log(3/\eps)$. Then let $E_2:\bits{3d_1}\times\bits{d}\to\bits{d_1}$ be an $(2d_1,\eps/3)$-extractor from Lemma~\ref{lemma:GUV}. Then we claim that $E:\bits{m+3d_1}\times\bits{d}\to\bits{m}$, defined as $E((x_1,x_2),s)=E_1(x_1,E_2(x_2,s))$, is a $(m+3d_1-\Delta,\eps)$ extractor, and hence a $(\eps,\delta)$ sampler by Lemma~\ref{lemma:Zuc}. 

To prove the claim, consider any $(m+3d_1-\Delta)$-source $X$. By Lemma~\ref{lemma:block}, $X$ is $(\eps/3)$-close to a $(m-\Delta,3d_1-\Delta-\log(3/\eps))$-block source $(X_1,X_2)\in\bits{3d_1}\times\bits{m}$. By Lemma~\ref{lemma:compo}, $E_1(X_1,E_2(X_2,U_d))$ is $2\eps/3$-close to uniform. Since $E(X,U_d)$ is $\eps/3$-close to $E_1(X_1,E_2(X_2,U_d))$, by triangle inequality it is $\eps$-close to uniform. Moreover, $d=O(\log (d_1/\eps))=O(\log\log(1/\delta)+\log(1/\eps))$, $n=m+3d_1=m+O(\log(1/\delta)+\log(1/\eps))$, and the required space to compute $E$ is $O(m+d_1+\log(1/\eps))=O(m+\log(1/\eps)+\log(1/\delta))$.
\end{proof}
